% Siavoosh Payandeh Azad Jan. 2019
\ProvidesFile{gradient.tex}[Cell background highlighting based on user data]
\RequirePackage{etoolbox}
\RequirePackage{pgf} % for calculating the values for gradient
\RequirePackage{xcolor} % enables the use of cellcolor make sure you have [table] option in the document class

% We enable a smooth color coding from red to yellow to green.
% 002400 ,
\definecolor{high}{HTML}{739B7F} % the color for the highest number in your data set (green)
\definecolor{medium}{HTML}{FFFFFF}  % The color in the medium (yellow)
\definecolor{low}{HTML}{FFFFFF}  % the color for the lowest number in your data set (red)

\newcommand*{\opacity}{90}% here you can change the opacity of the background color!

\definecolor{EHR_color}{HTML}{99A9CF}
\definecolor{WAVE_color}{HTML}{BFDAF1}
\definecolor{TVI_color}{HTML}{ECD9E8}

%======================================
% Data set related!
\newcommand*{\maxval}{0.9}% define the maximum value in your data set!
\newcommand*{\midval}{0.6}% define the minimum value on your data set
\newcommand*{\minval}{0.0}% define the minimum value on your data set
%======================================
% gradient function!
\newcommand{\gradient}[2]{
  % The values are calculated linearly between \minval \midval and \maxval
  \ifdimcomp{#1pt}{<}{\midval pt}{
    \pgfmathparse{int(round(100*(#1/(\midval-\minval))-(\minval*(100/(\midval-\minval)))))}
    \xdef\tempa{\pgfmathresult}
    \cellcolor{medium!\tempa!low!\opacity} #2
  }{
    \pgfmathparse{int(round(100*(#1/(\maxval-\midval))-(\midval*(100/(\maxval-\midval)))))}
    \xdef\tempa{\pgfmathresult}
    \cellcolor{high!\tempa!medium!\opacity} #2
  }
}
